\begin{proposition}[跃迁电流的连续性与规范端点项]
设 $\Psi_i,\Psi_f$ 都满足同一实外势下的 Dirac 方程。定义复数值跃迁源
\begin{equation*}
 \rho_{fi}^{\rm ch}=q_{\mathrm e}\Psi_f^\dagger\Psi_i,
 \qquad \vb{J}_{fi}^{\rm ch}=q_{\mathrm e}c\Psi_f^\dagger\boldsymbol\alpha\Psi_i.
\end{equation*}
它们满足
\begin{equation}
 \partial_t\rho_{fi}^{\rm ch}+\nabla\cdot\vb{J}_{fi}^{\rm ch}=0.
 \label{eq:feqo-transition-current-continuity}
\end{equation}
若空间边界项消失，则线性耦合矩阵元
$V_{fi}=\int(\rho_{fi}^{\rm ch}\phi-\vb{J}_{fi}^{\rm ch}\cdot\vb{A})\dd^3r$
在 $\vb{A}' =\vb{A}+\nabla\chi$、$\phi'=\phi-\partial_t\chi$ 下满足
\begin{equation}
 \delta V_{fi}(t)=-\frac{\dd}{\dd t}\int\dd^3r\,\rho_{fi}^{\rm ch}(\vb{r},t)\chi(\vb{r},t).
 \label{eq:feqo-transition-gauge-endpoint}
\end{equation}
完整跃迁概率保持规范不变；有限时间散射必须连同初末态的规范相位处理。
\end{proposition}

\begin{proof}
分别将 $\Psi_i$ 的哈密顿量方程左乘 $\Psi_f^\dagger$，以及 $\Psi_f$ 方程的厄米共轭右乘 $\Psi_i$。相加得到
\begin{align*}
 \partial_t(\Psi_f^\dagger\Psi_i)
 &=-c\big[(\partial_j\Psi_f^\dagger)\alpha_j\Psi_i
 +\Psi_f^\dagger\alpha_j\partial_j\Psi_i\big]\\
 &=-\nabla\cdot(c\Psi_f^\dagger\boldsymbol\alpha\Psi_i),
\end{align*}
因为质量矩阵和同一外势的两个局域乘法项逐项抵消。再对规范变化的矩阵元作空间分部积分：
\begin{align*}
 \delta V_{fi}
 &=-\int\dd^3r\,[\rho_{fi}^{\rm ch}\partial_t\chi
 +\vb{J}_{fi}^{\rm ch}\cdot\nabla\chi]\\
 &=-\int\dd^3r\,[\rho_{fi}^{\rm ch}\partial_t\chi
 +(\partial_t\rho_{fi}^{\rm ch})\chi].
\end{align*}
这正是总时间导数。自由定态的时间因子为 $\ee^{-\ii\omega_{if}t}$，$\omega_{if}=(E_i-E_f)/\hbar$，故频率分量上 $\nabla\cdot\vb{J}_{fi}^{\rm ch}=\ii\omega_{if}\rho_{fi}^{\rm ch}$。纵向跃迁电流与电荷密度必须满足这一约束。
\end{proof}

\begin{exbox}[正负能隙与有限脉冲的混合振幅]
只考察一个已指定的正负能通道，令自由能量差 $\Delta=E_f-E_i$，$|\Delta|\ge2m_{\mathrm e}c^2$，外势矩阵元为 $v_{fi}(t)$。在有限脉冲和一阶微扰条件下，求正负能通道混合振幅的界，并判断大能隙是否足以压低混合。
\textbf{解。} 在 $\eta=\hbar^{-1}\int\|V(t)\|\dd t\ll1$ 的有限时间微扰区间，一阶混合振幅为
\begin{equation*}
 c_{fi}^{(1)}=-\frac{\ii}{\hbar}\int_{t_i}^{t_f}\dd t\,
 v_{fi}(t)\ee^{\ii\Delta t/\hbar},
 \qquad c_{fi}=c_{fi}^{(1)}+O(\eta^2).
\end{equation*}
假定 $v_{fi}$ 绝对连续，分部积分给
\begin{equation*}
 |c_{fi}^{(1)}|\le
 \frac{|v_{fi}(t_i)|+|v_{fi}(t_f)|
 +\int_{t_i}^{t_f}|\dot v_{fi}(t)|\dd t}{|\Delta|}.
\end{equation*}
对于端点关闭、$v_{fi}(t)=v_*w(t/T)$ 的光滑脉冲，分子变为 $|v_*|\int|w'(s)|\dd s$。如果谱在能隙附近共振，$\dot v_{fi}$ 则包含该快速振荡，单纯比较峰值 $|v_*|$ 与能隙不再够用。低光子能量、外场在 Compton 尺度上缓变、有限脉冲混合小，是不同的控制条件；强场多光子与隧穿对产生不能仅由 $\hbar\omega<2m_{\mathrm e}c^2$ 排除。
\end{exbox}
